\documentclass{article}
\usepackage[portuguese]{babel}
\usepackage{listings}
\usepackage{xcolor}
\usepackage[margin=1.5cm]{geometry}

\definecolor{codegreen}{rgb}{0,0.6,0}
\definecolor{codegray}{rgb}{0.5,0.5,0.5}
\definecolor{codepurple}{rgb}{0.8,0,0.2}
\definecolor{backcolour}{rgb}{.8,.8,1}

\lstdefinestyle{mystyle}{
    backgroundcolor=\color{backcolour},   
    commentstyle=\color{codegreen},
    keywordstyle=\color{blue},
    numberstyle=\tiny\color{codegray},
    stringstyle=\color{codepurple},
    basicstyle=\ttfamily\footnotesize,
    breakatwhitespace=false,         
    breaklines=true,                 
    captionpos=b,                    
    keepspaces=true,                 
    numbers=left,                    
    numbersep=5pt,                  
    showspaces=false,                
    showstringspaces=false,
    showtabs=false,                  
    tabsize=2
}

\lstset{style=mystyle}
\begin{document}

\section{AA10}
Escreva um programa, em linguagem Python, para resolver um sistema de equações lineares triangular inferior a partir das seguintes etapas:
\begin{enumerate}
    \item ler a dimensão n da matriz dos coeficientes no console;
    \item registrar as entadas da matriz dos coeficientes no console;
    \item registrar as entradas do vetor dos resultados no console;
    \item imprimir o sistema original na forma de matriz ampliada;
    \item resolver e imprimir a solução do sistema de equações;
\end{enumerate}

\subsection{Solucao para sistema triangular superior}
O codigo para esse exercicio e praticamente igual ao do exercicio passado. Porem como no exercicio passado tivemos que trabalhar de traz pra frente, o codigo ficou mais dificil de entender. Dessa forma vamos explicar melhor aqui onde e mais facil de entender e dessa forma as alteracoes feitas no exercicio passado ficarao claras.

Por ser um sistema triangular superior. Obtemos o valor do primeiro elemento $x_1$ do vetor de incognitas $x$ simplesmente com 
\[
x_1 = \frac{b_1}{A_{[1,1]}}
\]
Tendo o valor de $x_1$, o utilizamos para encontrar $x_2$, usamos $x_1$ e $x_2$ para encontrar $x_3$ e continuamos assim ate finalizar o sistema usando:
\[
x_n =\frac{1}{A_{[n,n]}}[ b_n -\sum_{i=1}^{n} A_{[n,i]} \cdot x_i]
\]
\begin{lstlisting}[language=Python]
incognitas = [0]*ordem
incognitas[0] = vetorResultado[0]/matriz[0][0]
for n in range(1,ordem):
    somatorio = 0
    for i in range(n):
        somatorio += matriz[n][i] * incognitas[i]
    incognitas[n] = (vetorResultado[n]-somatorio)/matriz[n][n]
print(f"Solucao:{incognitas}")
\end{lstlisting}
Observe que temos que adaptar nossos indices devido a linguagem python utilizar 0 como primeiro elemento.
\subsection{Codigo completo:}
\begin{lstlisting}[language=Python]
deveParar = ""

def printmatriz(matriz):
    for linha in matriz:
        print(linha)

def printAmpliada(matriz, vetor):
    for i in range(len(matriz)):
        print(f"{matriz[i]}{vetor[i]}")


while (deveParar == ""):
    
    ordem = int(input("Digite a dimensao da matriz A\n"))
    matriz = []
    for u in range(ordem):
        linha = []
        for v in range(ordem):
            linha.append(int(input(f"Digite o valor na posisao {u+1},{v+1} \n")))
        matriz.append(linha)
    print("Matriz inicial A:")
    printmatriz(matriz)

    vetorResultado = []
    for i in range(ordem):
        vetorResultado.append(int(input(f"Digite o valor de X{i} do vetor resultado:\n")))
    printmatriz(vetorResultado)

    print("------00-----")
    print("Matriz Ampliada:")
    printAmpliada(matriz,vetorResultado)

    incognitas = [0]*ordem
    incognitas[0] = vetorResultado[0]/matriz[0][0]

    for n in range(1,ordem):
        somatorio = 0
        for i in range(n):
            somatorio += matriz[n][i] * incognitas[i]
        incognitas[n] = (vetorResultado[n]-somatorio)/matriz[n][n]

    print("Solucao:")
    print(incognitas)
    deveParar = input("Deseja continuar?(enter para continuar, qualquer tecla para parar)\n")

print("Ate mais")
\end{lstlisting}

\end{document}
